\problemname{Train Mania}
It is very well known that Julia likes the TV series \textit{På spåret}. So, it may not come as
a surprise that she enjoys traveling by train. Unfortunately, she has recently gone a bit off
track, and this once innocent hobby has turned into a manic need to ride trains. The problem is
that it is difficult to travel by train as much as she would like since trips take a long time.  

Even though she cannot take every train, she tries to travel with as many as possible. After
spending hours in her room trying to create the best travel plan, you decide to help her.
While the number of trains she takes is not exactly healthy, you prefer that she actually
rides trains rather than just planning her train journeys.  

Julia lives in the city of Lönnköping, which has $N$ stations numbered from $1$ to $N$.
To travel between them, there are two modes of transportation: electric scooters and trains.
There are so many electric scooters in the city that she can start a scooter ride from any station.
However, she can only ride a scooter along the $N-1$ roads in the city. It is possible to travel
between any two stations using one or more of these roads, meaning the road network forms a \textit{tree}.
Traveling along each road with a scooter takes a certain number of seconds.  

She can also take part in the next $K$ train departures. Each train follows a strict schedule: it
departs from a specific station at a specific time and arrives at its destination after a certain number
of seconds. If Julia is at a station at the exact second a train departs, she can take it if she wants.
She will then stay on the train until it reaches its final destination (she cannot bring herself to
get off a train before she absolutely has to). If a train arrives at a station at the exact second
another train departs, she can make the transfer if she chooses.  

Given the road network and the next $K$ train departures, how many trains can Julia take if she plans
her journey optimally? Since the plan is set for some time in the future, she is unsure where she
will start, so she wants you to calculate the maximum number of trains she can take for each possible
starting station. You can assume that she always begins her journey at time $0$.  


\section*{Input}
The first line of input contains two integers, $N$ and $K$ ($1 \leq N, K \leq 10^5$), representing
the number of stations in Lönnköping and the number of train departures to consider.

The next $N-1$ lines describe the road network. Each of these lines contains the integers
$u, v, s$ ($1 \leq u, v \leq N$, $1 \leq s \leq 10^9$), meaning there is a road between stations
$u$ and $v$. It takes $s$ seconds to travel from $u$ to $v$ or from $v$ to $u$ using a scooter.

After that, there are $K$ lines describing the train departures. Each of these lines contains
the integers $u, v, t, d$ ($1 \leq u, v \leq N$, $u \neq v$, $1 \leq t, d \leq 10^9$), meaning
that at time $t$, a train departs from station $u$ to station $v$ and arrives at station $v$ at
time $t+d$. It is guaranteed that all $t$ values are unique.

\section*{Output}
Let $x_i$ denote the maximum number of trains Julia can take if she starts her journey at station $i$.
Print $x_1, x_2, \dots, x_N$ separated by spaces.

\section*{Scoring}
Your solution will be tested on a set of test groups, each worth a number of points. Each test group contains
a set of test cases. To get the points for a test group you need to solve all test cases in the test group.

\noindent
\begin{tabular}{| l | l | p{12cm} |}
  \hline
  \textbf{Group} & \textbf{Points} & \textbf{Constraints} \\ \hline
  $1$    & $10$       & $N, K \leq 50$, $t, d \leq 100$ for all $t,d$. \\ \hline
  $2$    & $5$        & $N, K \leq 50$ \\ \hline
  $3$    & $12$       & $N \leq 1000, K \leq 200$ \\ \hline
  $4$    & $13$       & $N, K \leq 2000$ \\ \hline
  $5$    & $10$       & $K \leq 2000$ and the smallest $t$ is larger than the sum of all $s$. \\ \hline
  $6$    & $2$        & $K \leq 20000$ \\ \hline
  $7$    & $23$       & There is a path between station $i$ and station $i+1$ for all $1 \leq i \leq N-1$. \\ \hline
  $8$    & $25$       & No additional constraints. \\ \hline
\end{tabular}

